\chapter*{Wstęp}
\addcontentsline{toc}{chapter}{Wstęp}
\markboth{Wstęp}{Wstęp}
\label{ch:wstep}

Symulacje zjawisk opisanych równaniami różniczkowymi cząstkowymi należą do najczęściej
spotykanych zadań obliczeniowych w~nauce i~inżynierii. Wśród nich szczególne miejsce zajmują
obliczenia \emph{stencilowe}, czyli takie, w~których nowa wartość w~każdym punkcie siatki jest
funkcją wartości z punktów sąsiadujących.

Trójwymiarowe równanie przewodnictwa ciepła jest tradycyjnym przykładem takiego problemu.
Jego jawna dyskretyzacja różnicowa posiada niską intensywność arytmetyczną
oraz konieczność wymiany danych brzegowych między częściami dziedziny w~każdym kroku
czasowym~\cite{datta-stencil} co sprawia, że obliczenia stencilowe
są trudne do zrównoleglenia. Program symulujący to równanie nadaje się dobrze do badania granic skalowalności, ponieważ jego prosty wzorzec komunikacji, ograniczony do wymiany z~najbliższymi sąsiadami,
przy licznych węzłach i~tak obciąża współdzielone łącze sieciowe,
co ujawnia ograniczenia komunikacyjne klastra.

Programowanie równoległe tradycyjnie realizuje się za pomocą interfejsu
przekazywania komunikatów MPI (ang.~\emph{Message Passing Interface}), zyskując wysoką wydajność kosztem znacznej złożoności kodu, gdyż
programista jawnie zarządza podziałem danych, buforami wymiany oraz synchronizacją~\cite{mpi-standard}. Alternatywę
oferują języki i~biblioteki oparte na modelu \emph{PGAS} (ang.\ \emph{Partitioned Global Address
Space})~\cite{pgas}. Język Chapel,
rozwijany pierwotnie przez firmę Cray, jest
jednym z~bardziej rozwiniętych języków w~tej grupie~\cite{chapel-overview}. Umożliwia zapisywanie obliczeń
rozproszonych w~notacji zbliżonej do sekwencyjnej, z~globalnym widokiem tablic, wysokopoziomowymi
iteratorami równoległymi i~gotowymi rozkładami dziedzin. Szczegóły komunikacji pozostawia
warstwie wykonawczej, opartej na bibliotece GASNet.

Zbliżoną tematykę podjęto w~pracy współautorskiej
poświęconej równoległej symulacji dwuwymiarowego przewodnictwa ciepła na~klastrze Raspberry~Pi
oraz porównaniu implementacji w~językach Chapel, C\# i~Julia~\cite{falkowski2026parallel}, podczas
gdy niniejsza praca obejmuje przypadek trójwymiarowy na~klastrze Pionier.

Celem niniejszej pracy było zaprojektowanie, zaimplementowanie i~wdrożenie na~klastrze
programu symulującego równanie przewodnictwa ciepła w~Chapel oraz
analiza jego poprawności i~skalowalności.
